\chapter{Forms}

Forms drive most frontend/backend interaction: login, registration, and
every create/edit screen for tasks. This chapter covers the shared
pattern --- controlled inputs backed by one state object, validated
before submit, wired to the API with loading/error handling.

\section{Controlled Inputs and a Single Form State Object}

A controlled input's value is driven by React state, updated on every
keystroke via \texttt{onChange}. Keep the whole form as one object and use
a single generic handler keyed by each input's \texttt{name}.

\begin{lstlisting}[language=JavaScript]
const [form, setForm] = useState({ email: "", password: "" });

const handleChange = (e) => {
  setForm({ ...form, [e.target.name]: e.target.value });
};
// <input name="email" value={form.email} onChange={handleChange} />
\end{lstlisting}

\begin{notebox}[NOTE]
This one handler works for any number of fields as long as each
\texttt{name} matches a key in the form object --- add a field by adding a
key and an input, no new handler needed.
\end{notebox}

\section{Example 1: LoginForm}

\begin{lstlisting}[language=JavaScript]
// components/LoginForm.jsx
import { useState } from "react";
import api from "../api/axios";

function LoginForm({ onSuccess }) {
  const [form, setForm] = useState({ email: "", password: "" });
  const [errors, setErrors] = useState({});
  const [loading, setLoading] = useState(false);

  const handleChange = (e) => setForm({ ...form, [e.target.name]: e.target.value });

  const validate = () => {
    const errs = {};
    if (!form.email) errs.email = "Email is required";
    if (!form.password) errs.password = "Password is required";
    return errs;
  };

  const handleSubmit = async (e) => {
    e.preventDefault();
    const errs = validate();
    setErrors(errs);
    if (Object.keys(errs).length) return;
    setLoading(true);
    try {
      const res = await api.post("/auth/login", form);
      onSuccess(res.data.data);
    } catch (err) {
      setErrors({ form: err.response?.data?.message || "Login failed" });
    } finally {
      setLoading(false);
    }
  };

  return (
    <form onSubmit={handleSubmit} className="space-y-4">
      <input name="email" value={form.email} onChange={handleChange} placeholder="Email" />
      {errors.email && <p className="text-red-600 text-sm">{errors.email}</p>}
      <input name="password" type="password" value={form.password} onChange={handleChange} placeholder="Password" />
      {errors.password && <p className="text-red-600 text-sm">{errors.password}</p>}
      {errors.form && <p className="text-red-600 text-sm">{errors.form}</p>}
      <button type="submit" disabled={loading}>{loading ? "Logging in..." : "Log In"}</button>
    </form>
  );
}

export default LoginForm;
\end{lstlisting}

\section{Example 2: RegisterForm with Confirm Password}

Adds a \texttt{confirmPassword} field that's stripped before the request
is sent --- it exists only for client-side validation.

\begin{lstlisting}[language=JavaScript]
// components/RegisterForm.jsx (key differences from LoginForm)
const [form, setForm] = useState({ name: "", email: "", password: "", confirmPassword: "" });

const validate = () => {
  const errs = {};
  if (!form.name) errs.name = "Name is required";
  if (!/^\S+@\S+\.\S+$/.test(form.email)) errs.email = "Enter a valid email";
  if (form.password.length < 6) errs.password = "Password must be at least 6 characters";
  if (form.confirmPassword !== form.password) errs.confirmPassword = "Passwords do not match";
  return errs;
};

const handleSubmit = async (e) => {
  e.preventDefault();
  const errs = validate();
  setErrors(errs);
  if (Object.keys(errs).length) return;
  setLoading(true);
  try {
    const { confirmPassword, ...payload } = form; // never send this to the backend
    await api.post("/auth/register", payload);
  } catch (err) {
    setErrors({ form: err.response?.data?.message || "Registration failed" });
  } finally {
    setLoading(false);
  }
};
// Renders one input per key in `form` plus a matching error line, same shape as LoginForm.
\end{lstlisting}

\begin{warnbox}[WARNING]
Never send \texttt{confirmPassword} to the backend --- strip it via
destructuring before the API call, as shown above.
\end{warnbox}

\section{Example 3: TaskForm --- Create/Edit Dual Mode}

One form component handles both create and edit. An \texttt{id} prop picks
the mode: present means load-then-\texttt{PUT}, absent means blank-then-\texttt{POST}.

\begin{lstlisting}[language=JavaScript]
// components/TaskForm.jsx (delta from the pattern above)
function TaskForm({ id, onSaved }) {
  const isEditMode = Boolean(id);
  const [form, setForm] = useState({ title: "", description: "", status: "pending" });

  useEffect(() => {
    if (!isEditMode) return;
    api.get(`/tasks/${id}`).then((res) => {
      const { title, description, status } = res.data.data;
      setForm({ title, description, status });
    });
  }, [id, isEditMode]);

  const handleSubmit = async (e) => {
    e.preventDefault();
    if (!form.title.trim()) return setErrors({ title: "Title is required" });
    if (isEditMode) await api.put(`/tasks/${id}`, form);
    else await api.post("/tasks", form);
    onSaved();
  };
  // Inputs: title, description (textarea), status (select: pending/in-progress/done)
  // Submit label: isEditMode ? "Update Task" : "Create Task"
}
\end{lstlisting}

\begin{tipbox}[TIP]
This create/edit-by-prop pattern avoids two near-identical form
components: render \texttt{<TaskForm />} to create, \texttt{<TaskForm id={task.\_id} />} to edit.
\end{tipbox}

\section{Resetting Forms After Success}

Reset a create form to its initial values after success (see the create
branch above). Edit forms typically don't need a reset.

\whatwhyhow
{Forms use controlled inputs backed by a single state object, a generic \texttt{handleChange}, client-side validation before submit, and loading/error state around the API call.}
{Scattered per-field state makes validation and reset inconsistent; skipping client-side validation wastes a round trip.}
{Validate in a \texttt{validate()} call before the API request, set field errors from its return value, and reset on success.}
